\chapter{Ultradelayed Material Failure via Shear Banding After Straining in Networks} \label{sec:Ultradelayed}

\section{Introduction}
This chapter will study the behaviour of network materials under the step strain protocol outlined in \cref{sec:DeformationProtocols}. In particular, we will be exploring the same physics explored by Lockwood \textit{et al} in Ref. \cite{lockwood_ultradelayed_2025}. Using a variety of constitutive models, this work explored the formation of delayed dramatic shear banding instabilities long after the initial imposition of the strain. We hope that by introducing thermal effects to the simple network models described in \cref{sec:Networks}, we can find and explore how this effect is presented in a more realistic model. 

Shear banding and failure have been widely explored both experimentally and theoretically, as discussed in \cref{sec:ShearBanding}, and have shown to show failure in a range of protocols. Typically, protocols such as shear startup, outlined in \cref{sec:DeformationProtocols}, where at some time $t=0$ the system is started to be sheared with a rate $\dot{\gamma}$ have been used to study the formation of these bands \cite{nicolas_deformation_2018,moorcroft_shear_2014,moorcroft_criteria_2013}. 

\section{Methodology}
To explore the ultra delayed failure of networks, we extend the network model defined in \cref{sec:Networks}, in particular, taking the Jamming-based networks defined in \cref{sec:PackingBasedNetworks}, evolving them under the dynamics discussed in \cref{sec:NetworkDynamics} and introducing a thermal component to these models. 


\section{Brownian Dynamics}
As discussed in \cref{sec:Fracture}, the fracture of bonds in the network plays a large role in determining the material's response and properties. 

\section{Boltzman Fracture}
\section{Stress Decay and Strain Response}
\section{Banding Time $t^*$ and the Imposed Step Strain $\gamma_0$}
\section{Banding Time $t^*$ and Temperature $x$}

\section{Conclusion}